\documentclass[12pt,a4paper]{article}

\usepackage[T2A]{fontenc}
\usepackage[utf8]{inputenc}
\usepackage[russian]{babel}
\usepackage{amsmath,amssymb,amsthm}
\usepackage{mathtools}
\usepackage{enumitem}
\usepackage{listings}
\usepackage{xcolor}
\usepackage{array}
\usepackage{hyperref}
\usepackage{stmaryrd}

\lstdefinestyle{code}{
  basicstyle=\ttfamily\small,
  keywordstyle=\color{blue},
  commentstyle=\color{gray},
  stringstyle=\color{red!60!black},
  columns=fullflexible,
  keepspaces=true,
  frame=single,
  breaklines=true,
  showstringspaces=false
}

\lstset{style=code}

\newtheorem{definition}{Определение}
\newtheorem{theorem}{Теорема}
\newtheorem{lemma}{Лемма}
\newtheorem{example}{Пример}
\newtheorem{remark}{Замечание}

\title{Билет №31 \\ \small Языки программирования. Процедурные языки программирования (Fortran, C), Функциональные языки программирования (LISP, ML, Haskell), логическое программирование (Prolog), объектно-ориентированные языки программирования (Java, C\#). (СпБГУ) \\ Основные понятия логического программирования. Теорема Эрбрана. Метод резолюций. (СпБПУ)}
\author{Сериков Владислав Александрович}
\date{13 мая 2026}

\begin{document}

\maketitle
\tableofcontents
\newpage

\section{Языки программирования}

\subsection{Понятие языка программирования}

\begin{definition}
\textbf{Язык программирования} --- формальная система записи алгоритмов и структур данных, предназначенная для описания вычислений, выполняемых человеком или машиной.
\end{definition}

Язык программирования обычно задаётся следующими компонентами:
\begin{enumerate}
    \item \textbf{Алфавит} --- множество допустимых символов.
    \item \textbf{Лексика} --- правила образования токенов: идентификаторов, чисел, ключевых слов, операций.
    \item \textbf{Синтаксис} --- правила построения корректных программ.
    \item \textbf{Статическая семантика} --- правила корректности, проверяемые до выполнения программы, например правила типов и областей видимости.
    \item \textbf{Динамическая семантика} --- смысл выполнения программы.
    \item \textbf{Прагматика} --- практические аспекты использования языка: эффективность, переносимость, удобство сопровождения.
\end{enumerate}

\begin{definition}
\textbf{Парадигма программирования} --- совокупность идей, принципов и моделей построения программ.
\end{definition}

Основные парадигмы:
\begin{itemize}
    \item процедурное программирование;
    \item функциональное программирование;
    \item логическое программирование;
    \item объектно-ориентированное программирование;
    \item декларативное программирование;
    \item императивное программирование.
\end{itemize}

\subsection{Императивный и декларативный подходы}

\begin{definition}
\textbf{Императивное программирование} --- стиль программирования, при котором программа описывает последовательность команд, изменяющих состояние вычислительной системы.
\end{definition}

К императивным языкам относятся Fortran, C, Pascal, Ada, Java, C\# в их императивной части.

\begin{definition}
\textbf{Декларативное программирование} --- стиль программирования, при котором программа описывает, \emph{что} требуется вычислить, а не \emph{как именно} выполнять вычисления.
\end{definition}

К декларативным языкам относятся SQL, Prolog, Haskell, ML в значительной своей части.

\section{Процедурные языки программирования}

\subsection{Общая характеристика процедурного программирования}

\begin{definition}
\textbf{Процедурное программирование} --- разновидность императивного программирования, в которой программа строится как совокупность процедур и функций, последовательно вызывающих друг друга.
\end{definition}

Основные признаки процедурных языков:
\begin{itemize}
    \item наличие переменных и присваивания;
    \item последовательное выполнение команд;
    \item условные операторы;
    \item циклы;
    \item процедуры и функции;
    \item стек вызовов;
    \item локальные и глобальные области видимости.
\end{itemize}

Типичная структура процедурной программы:
$$
\text{данные} + \text{процедуры обработки данных}.
$$

\subsection{Модель вычислений процедурного языка}

Состояние программы можно представить как отображение
$$
\sigma : Var \to Val,
$$
где $Var$ --- множество переменных, $Val$ --- множество значений.

Оператор присваивания
$$
x := e
$$
изменяет состояние:
$$
\sigma' = \sigma[x \mapsto \llbracket e \rrbracket_\sigma],
$$
где $\llbracket e \rrbracket_\sigma$ --- значение выражения $e$ в состоянии $\sigma$.

\subsection{Fortran}

\subsubsection{Общая характеристика}

Fortran, от \emph{Formula Translation}, является одним из первых языков высокого уровня. Он создавался прежде всего для научных и инженерных вычислений.

Характерные особенности Fortran:
\begin{itemize}
    \item развитая поддержка численных вычислений;
    \item эффективная работа с массивами;
    \item наличие компилируемых программ;
    \item ориентация на высокую производительность;
    \item широкое применение в задачах моделирования, физики, вычислительной математики.
\end{itemize}

\subsubsection{Минимальный пример на Fortran}

\begin{lstlisting}[language=Fortran]
program sum_array
  implicit none
  integer :: i
  real :: a(5), s

  a = (/1.0, 2.0, 3.0, 4.0, 5.0/)
  s = 0.0

  do i = 1, 5
     s = s + a(i)
  end do

  print *, s
end program sum_array
\end{lstlisting}

Алгоритм суммирования массива:
\begin{enumerate}
    \item Инициализировать сумму $s := 0$.
    \item Для каждого элемента массива $a_i$ выполнить $s := s + a_i$.
    \item После обработки всех элементов вывести $s$.
\end{enumerate}

Для массива
$$
a = [1,2,3,4,5]
$$
шаги выглядят так:
$$
0 \to 1 \to 3 \to 6 \to 10 \to 15.
$$

\subsection{Язык C}

\subsubsection{Общая характеристика}

C --- процедурный язык системного программирования. Он сочетает высокоуровневые конструкции с возможностью низкоуровневого управления памятью.

Основные особенности C:
\begin{itemize}
    \item статическая типизация;
    \item указатели;
    \item ручное управление памятью;
    \item компактный синтаксис;
    \item прямая связь с моделью памяти машины;
    \item широкое применение в операционных системах, компиляторах, встраиваемых системах.
\end{itemize}

\subsubsection{Минимальный пример на C}

\begin{lstlisting}[language=C]
#include <stdio.h>

int sum(int *a, int n) {
    int s = 0;
    for (int i = 0; i < n; i++) {
        s += a[i];
    }
    return s;
}

int main(void) {
    int a[] = {1, 2, 3, 4, 5};
    printf("%d\n", sum(a, 5));
    return 0;
}
\end{lstlisting}

\subsubsection{Указатели}

\begin{definition}
\textbf{Указатель} --- переменная, значением которой является адрес некоторого объекта в памяти.
\end{definition}

Пример:
\begin{lstlisting}[language=C]
int x = 10;
int *p = &x;
*p = 20;
\end{lstlisting}

После выполнения кода значение $x$ равно $20$.

\subsection{Достоинства и недостатки процедурного программирования}

Достоинства:
\begin{itemize}
    \item естественное описание пошаговых алгоритмов;
    \item высокая эффективность;
    \item близость к архитектуре фон Неймана;
    \item хорошая поддержка компиляторами.
\end{itemize}

Недостатки:
\begin{itemize}
    \item сложность управления изменяемым состоянием;
    \item возможное появление побочных эффектов;
    \item трудность сопровождения больших программ без дополнительной модульной организации;
    \item ошибки при ручном управлении памятью.
\end{itemize}

\section{Функциональные языки программирования}

\subsection{Основные идеи функционального программирования}

\begin{definition}
\textbf{Функциональное программирование} --- парадигма программирования, в которой вычисление рассматривается как вычисление значений функций, а программа строится из выражений и функций.
\end{definition}

Основные идеи:
\begin{itemize}
    \item функции являются значениями первого класса;
    \item широкое использование рекурсии;
    \item неизменяемость данных;
    \item отсутствие или ограничение побочных эффектов;
    \item функции высшего порядка;
    \item композиционность;
    \item ленивые вычисления в некоторых языках.
\end{itemize}

\begin{definition}
\textbf{Функция первого класса} --- функция, которую можно передавать как аргумент, возвращать как результат и присваивать переменной.
\end{definition}

\begin{definition}
\textbf{Функция высшего порядка} --- функция, принимающая функцию в качестве аргумента или возвращающая функцию как результат.
\end{definition}

\subsection{Лямбда-исчисление как основа функционального программирования}

\begin{definition}
\textbf{Лямбда-исчисление} --- формальная система, предназначенная для описания функций, их применения и вычислений.
\end{definition}

Термы чистого лямбда-исчисления:
$$
M,N ::= x \mid (\lambda x.M) \mid (MN).
$$

Здесь:
\begin{itemize}
    \item $x$ --- переменная;
    \item $\lambda x.M$ --- абстракция, то есть функция с параметром $x$ и телом $M$;
    \item $MN$ --- применение функции $M$ к аргументу $N$.
\end{itemize}

Основное правило вычисления --- $\beta$-редукция:
$$
(\lambda x.M)N \to_\beta M[x := N],
$$
где $M[x := N]$ означает подстановку $N$ вместо свободных вхождений $x$ в $M$.

\begin{example}
$$
(\lambda x. x + 1)\ 5 \to 5 + 1 \to 6.
$$
\end{example}

\subsection{LISP}

\subsubsection{Общая характеристика}

LISP --- один из старейших функциональных языков. Его название происходит от \emph{LISt Processing}. Язык особенно важен для истории искусственного интеллекта и символьных вычислений.

Особенности LISP:
\begin{itemize}
    \item программы и данные имеют сходное представление в виде списков;
    \item активное использование рекурсии;
    \item динамическая типизация;
    \item интерактивная среда разработки;
    \item макросы и метапрограммирование.
\end{itemize}

\subsubsection{S-выражения}

Основная форма записи в LISP --- S-выражение:
$$
(f\ a_1\ a_2\ \dots\ a_n).
$$

Пример:
\begin{lstlisting}[language=Lisp]
(+ 1 2 3)
\end{lstlisting}

Результат:
$$
6.
$$

\subsubsection{Рекурсивная функция на LISP}

\begin{lstlisting}[language=Lisp]
(defun fact (n)
  (if (= n 0)
      1
      (* n (fact (- n 1)))))
\end{lstlisting}

Вычисление:
$$
fact(4) = 4 \cdot fact(3)
        = 4 \cdot 3 \cdot fact(2)
        = 4 \cdot 3 \cdot 2 \cdot fact(1)
        = 4 \cdot 3 \cdot 2 \cdot 1 \cdot fact(0)
        = 24.
$$

\subsection{ML}

\subsubsection{Общая характеристика}

ML --- семейство функциональных языков со строгой статической типизацией и выводом типов.

К важнейшим особенностям ML относятся:
\begin{itemize}
    \item статическая типизация;
    \item параметрический полиморфизм;
    \item автоматический вывод типов;
    \item сопоставление с образцом;
    \item алгебраические типы данных;
    \item модули.
\end{itemize}

\subsubsection{Пример на Standard ML}

\begin{lstlisting}
fun fact 0 = 1
  | fact n = n * fact (n - 1);
\end{lstlisting}

Тип функции:
$$
fact : int \to int.
$$

\subsubsection{Сопоставление с образцом}

\begin{lstlisting}
fun length [] = 0
  | length (_::xs) = 1 + length xs;
\end{lstlisting}

Для списка $[a,b,c]$:
$$
length([a,b,c]) = 1 + length([b,c])
= 1 + 1 + length([c])
= 1 + 1 + 1 + length([])
= 3.
$$

\subsection{Haskell}

\subsubsection{Общая характеристика}

Haskell --- чистый функциональный язык программирования с ленивыми вычислениями и мощной статической системой типов.

Особенности Haskell:
\begin{itemize}
    \item чистота функций;
    \item неизменяемость данных;
    \item ленивые вычисления;
    \item алгебраические типы данных;
    \item классы типов;
    \item монады для контролируемого описания эффектов.
\end{itemize}

\subsubsection{Минимальный пример на Haskell}

\begin{lstlisting}[language=Haskell]
fact :: Integer -> Integer
fact 0 = 1
fact n = n * fact (n - 1)
\end{lstlisting}

\subsubsection{Функции высшего порядка}

\begin{lstlisting}[language=Haskell]
map (\x -> x * x) [1,2,3,4]
\end{lstlisting}

Результат:
$$
[1,4,9,16].
$$

Алгоритм работы \texttt{map f xs}:
\begin{enumerate}
    \item Если список пуст, вернуть пустой список.
    \item Иначе применить $f$ к голове списка.
    \item Рекурсивно применить \texttt{map f} к хвосту списка.
    \item Соединить результат.
\end{enumerate}

Для
$$
f(x)=x^2,\quad xs=[1,2,3]
$$
получаем:
$$
map(f,[1,2,3])
=
[f(1),f(2),f(3)]
=
[1,4,9].
$$

\subsection{Преимущества и недостатки функционального программирования}

Преимущества:
\begin{itemize}
    \item близость к математическому стилю рассуждений;
    \item удобство доказательства корректности;
    \item отсутствие многих ошибок, связанных с изменяемым состоянием;
    \item хорошая пригодность для параллельных вычислений;
    \item выразительные абстракции.
\end{itemize}

Недостатки:
\begin{itemize}
    \item непривычность для программистов, воспитанных на императивном стиле;
    \item потенциальные сложности с эффективностью;
    \item необходимость специальных механизмов для ввода-вывода и эффектов;
    \item иногда менее очевидное управление памятью и временем выполнения.
\end{itemize}

\section{Логическое программирование}

\subsection{Общая характеристика}

\begin{definition}
\textbf{Логическое программирование} --- парадигма программирования, в которой программа задаётся как множество логических формул, а вычисление состоит в логическом выводе следствий из этих формул.
\end{definition}

В процедурной программе задаётся последовательность действий. В логической программе задаются факты и правила, а система вывода ищет доказательство цели.

Общая идея:
$$
\text{программа} + \text{запрос} \Longrightarrow \text{логический вывод ответа}.
$$

\subsection{Prolog}

\subsubsection{Основные элементы Prolog}

Prolog-программа состоит из:
\begin{itemize}
    \item фактов;
    \item правил;
    \item запросов.
\end{itemize}

Факт:
\begin{lstlisting}[language=Prolog]
parent(alice, bob).
\end{lstlisting}

Правило:
\begin{lstlisting}[language=Prolog]
grandparent(X, Z) :- parent(X, Y), parent(Y, Z).
\end{lstlisting}

Запрос:
\begin{lstlisting}[language=Prolog]
?- grandparent(alice, Z).
\end{lstlisting}

\subsubsection{Смысл правила}

Правило
\begin{lstlisting}[language=Prolog]
grandparent(X, Z) :- parent(X, Y), parent(Y, Z).
\end{lstlisting}
можно прочитать как:
$$
grandparent(X,Z) \leftarrow parent(X,Y) \land parent(Y,Z).
$$

Или в логической форме:
$$
\forall X \forall Y \forall Z
\bigl(parent(X,Y) \land parent(Y,Z) \to grandparent(X,Z)\bigr).
$$

\subsection{Унификация}

\begin{definition}
\textbf{Унификация} --- процесс нахождения подстановки, которая делает два терма синтаксически одинаковыми.
\end{definition}

\begin{definition}
\textbf{Подстановка} --- конечное отображение переменных в термы:
$$
\theta = \{x_1 \mapsto t_1,\dots,x_n \mapsto t_n\}.
$$
\end{definition}

Если для термов $s$ и $t$ существует подстановка $\theta$ такая, что
$$
s\theta = t\theta,
$$
то термы называются унифицируемыми.

\begin{definition}
\textbf{Наиболее общий унификатор} двух термов $s$ и $t$ --- такая подстановка $\theta$, что:
$$
s\theta = t\theta,
$$
и для любого другого унификатора $\sigma$ существует подстановка $\delta$, для которой
$$
\sigma = \theta \delta.
$$
\end{definition}

\subsubsection{Алгоритм унификации}

Пусть требуется унифицировать множество уравнений:
$$
E = \{s_1 = t_1,\dots,s_n=t_n\}.
$$

Алгоритм:
\begin{enumerate}
    \item Если $E$ пусто, вернуть пустую подстановку.
    \item Выбрать уравнение $s=t$ из $E$.
    \item Если $s$ и $t$ одинаковы, удалить уравнение.
    \item Если $s$ --- переменная $x$, а $x$ не входит в $t$, заменить $x$ на $t$ во всех остальных уравнениях и добавить подстановку $x \mapsto t$.
    \item Если $t$ --- переменная $x$, а $x$ не входит в $s$, заменить $x$ на $s$ во всех остальных уравнениях и добавить подстановку $x \mapsto s$.
    \item Если $s=f(s_1,\dots,s_n)$ и $t=f(t_1,\dots,t_n)$, заменить уравнение $s=t$ на систему
    $$
    s_1=t_1,\dots,s_n=t_n.
    $$
    \item Если главные функциональные символы различны или число аргументов различно, унификация невозможна.
    \item Если переменная должна быть заменена термом, содержащим эту же переменную, например $x=f(x)$, унификация невозможна. Это называется проверкой вхождения.
\end{enumerate}

\subsubsection{Пример унификации}

Унифицировать:
$$
p(X,f(Y),Y)
\quad\text{и}\quad
p(a,f(b),Z).
$$

Получаем систему:
$$
X=a,\quad f(Y)=f(b),\quad Y=Z.
$$

Из $f(Y)=f(b)$ следует:
$$
Y=b.
$$

Из $Y=Z$ и $Y=b$ следует:
$$
Z=b.
$$

Итак,
$$
\theta = \{X \mapsto a,\; Y \mapsto b,\; Z \mapsto b\}.
$$

\section{Объектно-ориентированные языки программирования}

\subsection{Основные понятия}

\begin{definition}
\textbf{Объектно-ориентированное программирование} --- парадигма, в которой программа строится как совокупность взаимодействующих объектов, объединяющих состояние и поведение.
\end{definition}

\begin{definition}
\textbf{Объект} --- сущность, обладающая состоянием, поведением и идентичностью.
\end{definition}

\begin{definition}
\textbf{Класс} --- описание структуры и поведения объектов некоторого вида.
\end{definition}

Ключевые принципы ООП:
\begin{enumerate}
    \item инкапсуляция;
    \item наследование;
    \item полиморфизм;
    \item абстракция.
\end{enumerate}

\subsection{Инкапсуляция}

\begin{definition}
\textbf{Инкапсуляция} --- объединение данных и методов, работающих с этими данными, с сокрытием внутреннего представления объекта.
\end{definition}

Пример на Java:
\begin{lstlisting}[language=Java]
class Counter {
    private int value = 0;

    public void increment() {
        value++;
    }

    public int getValue() {
        return value;
    }
}
\end{lstlisting}

Поле \texttt{value} скрыто от внешнего кода, доступ к нему осуществляется через методы.

\subsection{Наследование}

\begin{definition}
\textbf{Наследование} --- механизм создания нового класса на основе существующего с повторным использованием и расширением его свойств и методов.
\end{definition}

Пример:
\begin{lstlisting}[language=Java]
class Animal {
    void speak() {
        System.out.println("Some sound");
    }
}

class Dog extends Animal {
    @Override
    void speak() {
        System.out.println("Woof");
    }
}
\end{lstlisting}

\subsection{Полиморфизм}

\begin{definition}
\textbf{Полиморфизм} --- возможность использовать объекты разных классов через общий интерфейс.
\end{definition}

Пример:
\begin{lstlisting}[language=Java]
Animal a = new Dog();
a.speak();
\end{lstlisting}

Хотя переменная имеет тип \texttt{Animal}, вызывается метод класса \texttt{Dog}. Это пример динамического связывания.

\subsection{Java}

\subsubsection{Общая характеристика}

Java --- объектно-ориентированный язык программирования со статической типизацией и управляемой средой выполнения.

Особенности Java:
\begin{itemize}
    \item компиляция в байт-код;
    \item выполнение на виртуальной машине JVM;
    \item автоматическое управление памятью;
    \item сборка мусора;
    \item развитая стандартная библиотека;
    \item отсутствие множественного наследования классов;
    \item интерфейсы;
    \item исключения;
    \item многопоточность.
\end{itemize}

Минимальный пример:
\begin{lstlisting}[language=Java]
class Hello {
    public static void main(String[] args) {
        System.out.println("Hello");
    }
}
\end{lstlisting}

\subsection{C\#}

\subsubsection{Общая характеристика}

C\# --- объектно-ориентированный язык программирования платформы .NET. Он сочетает идеи C++, Java и современных языков с развитой системой типов.

Особенности C\#:
\begin{itemize}
    \item статическая типизация;
    \item выполнение в среде CLR;
    \item автоматическая сборка мусора;
    \item свойства;
    \item делегаты;
    \item события;
    \item LINQ;
    \item асинхронное программирование через \texttt{async}/\texttt{await}.
\end{itemize}

Минимальный пример:
\begin{lstlisting}[language={[Sharp]C}]
using System;

class Program {
    static void Main() {
        Console.WriteLine("Hello");
    }
}
\end{lstlisting}

\subsection{Сравнение парадигм}

\begin{center}
\begin{tabular}{|m{3cm}|m{4cm}|m{5cm}|}
\hline
\textbf{Парадигма} & \textbf{Основная идея} & \textbf{Примеры языков} \\
\hline
Процедурная & Последовательность команд и процедур & Fortran, C, Pascal \\
\hline
Функциональная & Вычисление значений функций & LISP, ML, Haskell \\
\hline
Логическая & Вывод следствий из фактов и правил & Prolog \\
\hline
Объектно-ориентированная & Взаимодействие объектов & Java, C\#, C++ \\
\hline
\end{tabular}
\end{center}

\section{Основные понятия логического программирования}

\subsection{Синтаксис логики первого порядка}

Логическое программирование обычно основано на фрагменте логики первого порядка.

\begin{definition}
\textbf{Сигнатура} --- набор функциональных и предикатных символов с указанными арностями.
\end{definition}

Например, сигнатура может содержать:
$$
a,\ b,\ f/1,\ g/2,\ parent/2,\ less/2.
$$

\begin{definition}
\textbf{Терм} определяется рекурсивно:
\begin{enumerate}
    \item переменная является термом;
    \item константа является термом;
    \item если $f$ --- $n$-местный функциональный символ, а $t_1,\dots,t_n$ --- термы, то $f(t_1,\dots,t_n)$ --- терм.
\end{enumerate}
\end{definition}

\begin{definition}
\textbf{Атомарная формула} имеет вид
$$
P(t_1,\dots,t_n),
$$
где $P$ --- $n$-местный предикатный символ, а $t_1,\dots,t_n$ --- термы.
\end{definition}

Формулы строятся из атомов с помощью связок:
$$
\neg,\ \land,\ \lor,\ \to,\ \leftrightarrow
$$
и кванторов:
$$
\forall,\ \exists.
$$

\subsection{Литералы, дизъюнкты и клаузы}

\begin{definition}
\textbf{Литерал} --- атомарная формула или её отрицание.
\end{definition}

Примеры литералов:
$$
P(x),\quad \neg Q(f(a),y).
$$

\begin{definition}
\textbf{Дизъюнкт}, или \textbf{клауза}, --- дизъюнкция конечного числа литералов:
$$
L_1 \lor L_2 \lor \dots \lor L_n.
$$
\end{definition}

Пустая клауза обозначается
$$
\square
$$
и соответствует противоречию.

\subsection{Хорновские дизъюнкты}

\begin{definition}
\textbf{Хорновская клауза} --- дизъюнкт, содержащий не более одного положительного литерала.
\end{definition}

Например:
$$
\neg A_1 \lor \dots \lor \neg A_n \lor B.
$$

Эта формула эквивалентна правилу:
$$
A_1 \land \dots \land A_n \to B.
$$

В Prolog она записывается как:
\begin{lstlisting}[language=Prolog]
B :- A1, ..., An.
\end{lstlisting}

Частные случаи:
\begin{enumerate}
    \item Факт:
    $$
    B.
    $$
    В Prolog:
\begin{lstlisting}[language=Prolog]
B.
\end{lstlisting}

    \item Цель:
    $$
    \neg A_1 \lor \dots \lor \neg A_n.
    $$
    В Prolog:
\begin{lstlisting}[language=Prolog]
?- A1, ..., An.
\end{lstlisting}
\end{enumerate}

\subsection{Декларативная и процедурная семантика Prolog}

\subsubsection{Декларативная семантика}

Декларативно Prolog-программа задаёт множество логических следствий. Запрос успешен, если он логически следует из программы.

Если программа $P$ содержит факты и правила, а цель $G$, то ответ ищется из отношения:
$$
P \models G\theta,
$$
где $\theta$ --- подстановка для переменных запроса.

\subsubsection{Процедурная семантика}

Процедурно Prolog использует:
\begin{itemize}
    \item унификацию;
    \item поиск в глубину;
    \item возврат, или backtracking;
    \item стратегию выбора целей слева направо.
\end{itemize}

\subsection{SLD-резолюция}

\begin{definition}
\textbf{SLD-резолюция} --- специальный вариант резолюции для определённых хорновских программ, применяемый в Prolog.
\end{definition}

Пусть имеется цель:
$$
\leftarrow A_1,\dots,A_i,\dots,A_n
$$
и правило:
$$
B \leftarrow C_1,\dots,C_m.
$$

Если $A_i$ и $B$ унифицируются с наиболее общим унификатором $\theta$, то новая цель:
$$
\leftarrow (A_1,\dots,A_{i-1},C_1,\dots,C_m,A_{i+1},\dots,A_n)\theta.
$$

\subsubsection{Пример SLD-вывода}

Программа:
\begin{lstlisting}[language=Prolog]
parent(alice, bob).
parent(bob, charlie).

grandparent(X, Z) :- parent(X, Y), parent(Y, Z).
\end{lstlisting}

Запрос:
\begin{lstlisting}[language=Prolog]
?- grandparent(alice, Z).
\end{lstlisting}

Шаги вывода:
$$
\leftarrow grandparent(alice,Z).
$$

Используем правило:
$$
grandparent(X,Z_1) \leftarrow parent(X,Y), parent(Y,Z_1).
$$

Унификация:
$$
grandparent(alice,Z) = grandparent(X,Z_1)
$$
даёт
$$
\theta_1 = \{X \mapsto alice,\ Z_1 \mapsto Z\}.
$$

Новая цель:
$$
\leftarrow parent(alice,Y), parent(Y,Z).
$$

Первый подзапрос унифицируется с фактом:
$$
parent(alice,bob).
$$

Получаем:
$$
\theta_2 = \{Y \mapsto bob\}.
$$

Новая цель:
$$
\leftarrow parent(bob,Z).
$$

Она унифицируется с фактом:
$$
parent(bob,charlie).
$$

Получаем:
$$
\theta_3 = \{Z \mapsto charlie\}.
$$

Итоговый ответ:
$$
Z = charlie.
$$

\section{Теорема Эрбрана}

\subsection{Эрбранов универсум и эрбранов базис}

\begin{definition}
\textbf{Эрбранов универсум} $HU$ для сигнатуры --- множество всех замкнутых термов, которые можно построить из констант и функциональных символов данной сигнатуры.
\end{definition}

Если в сигнатуре нет констант, добавляют специальную константу, например $c$.

\begin{example}
Пусть сигнатура содержит константу $a$ и одноместный функциональный символ $f$.
Тогда:
$$
HU = \{a,\ f(a),\ f(f(a)),\ f(f(f(a))),\dots\}.
$$
\end{example}

\begin{definition}
\textbf{Эрбранов базис} $HB$ --- множество всех замкнутых атомов, которые можно построить из предикатных символов и термов эрбранова универсума.
\end{definition}

\begin{example}
Если имеется предикат $P/1$ и
$$
HU=\{a,f(a),f(f(a)),\dots\},
$$
то
$$
HB =
\{P(a), P(f(a)), P(f(f(a))),\dots\}.
$$
\end{example}

\begin{definition}
\textbf{Эрбранова интерпретация} --- интерпретация, областью которой является эрбранов универсум, а каждый функциональный символ интерпретируется сам собой.
\end{definition}

В эрбрановой интерпретации значение терма совпадает с самим термом.

\subsection{Основная идея теоремы Эрбрана}

Теорема Эрбрана связывает выполнимость формул первого порядка с выполнимостью множества их замкнутых подстановочных экземпляров.

Идея состоит в следующем:
\begin{itemize}
    \item формула первого порядка может содержать кванторы и переменные;
    \item вместо рассуждения обо всех структурах можно рассматривать замкнутые термы эрбранова универсума;
    \item невыполнимость множества дизъюнктов первого порядка эквивалентна невыполнимости некоторого конечного множества его замкнутых экземпляров.
\end{itemize}

\subsection{Формулировка теоремы Эрбрана для множеств дизъюнктов}

\begin{theorem}[Теорема Эрбрана]
Пусть $S$ --- множество дизъюнктов логики первого порядка, а $S^*$ --- множество всех замкнутых экземпляров дизъюнктов из $S$, полученных подстановкой термов из эрбранова универсума вместо переменных.

Тогда множество $S$ невыполнимо тогда и только тогда, когда существует конечное множество
$$
S_0 \subseteq S^*
$$
такое, что $S_0$ невыполнимо в смысле исчисления высказываний.
\end{theorem}

\subsection{Доказательство теоремы Эрбрана}

\begin{proof}
Докажем обе импликации.

\textbf{1. Если существует конечное невыполнимое множество замкнутых экземпляров, то $S$ невыполнимо.}

Пусть
$$
S_0 \subseteq S^*
$$
и $S_0$ невыполнимо в смысле логики высказываний.

Предположим противное: пусть $S$ выполнимо в некоторой интерпретации $\mathcal{I}$.

Каждый элемент $C' \in S_0$ является замкнутым экземпляром некоторого дизъюнкта $C \in S$. То есть
$$
C' = C\theta,
$$
где $\theta$ --- подстановка замкнутых термов вместо переменных.

Если интерпретация $\mathcal{I}$ удовлетворяет $C$, то она удовлетворяет и любой его замкнутый экземпляр $C\theta$. Следовательно, $\mathcal{I}$ удовлетворяет всем элементам $S_0$.

Тогда $S_0$ выполнимо, что противоречит предположению о его невыполнимости.

Значит, $S$ невыполнимо.

\medskip

\textbf{2. Если $S$ невыполнимо, то существует конечное невыполнимое множество замкнутых экземпляров.}

Докажем контрапозицию.

Предположим, что любое конечное подмножество $S^*$ выполнимо в смысле логики высказываний. Требуется показать, что $S$ выполнимо.

Так как каждое конечное подмножество $S^*$ выполнимо, по теореме компактности для логики высказываний всё множество $S^*$ выполнимо в некоторой булевой оценке $v$.

Булева оценка $v$ задаёт истинность каждого замкнутого атома из эрбранова базиса:
$$
v : HB \to \{0,1\}.
$$

Построим эрбранову интерпретацию $\mathcal{H}$:
\begin{enumerate}
    \item область интерпретации равна эрбранову универсуму $HU$;
    \item каждый функциональный символ интерпретируется как конструктор термов;
    \item каждый предикатный символ $P$ интерпретируется так:
    $$
    \mathcal{H} \models P(t_1,\dots,t_n)
    \quad\Longleftrightarrow\quad
    v(P(t_1,\dots,t_n)) = 1.
    $$
\end{enumerate}

Покажем, что $\mathcal{H}$ удовлетворяет $S$.

Возьмём произвольный дизъюнкт $C \in S$. Нужно показать, что он истинен при любой оценке переменных. В эрбрановой интерпретации любая оценка переменных назначает переменным элементы из $HU$, то есть замкнутые термы.

Пусть оценка переменных соответствует подстановке $\theta$ замкнутых термов. Тогда значение $C$ при этой оценке совпадает со значением замкнутого экземпляра
$$
C\theta \in S^*.
$$

Но все формулы из $S^*$ истинны при оценке $v$, а интерпретация $\mathcal{H}$ построена в точности по этой оценке. Значит,
$$
\mathcal{H} \models C\theta.
$$

Следовательно,
$$
\mathcal{H} \models C.
$$

Так как $C$ выбран произвольно, то
$$
\mathcal{H} \models S.
$$

Значит, $S$ выполнимо.

Контрапозиция доказана. Следовательно, если $S$ невыполнимо, то существует конечное невыполнимое множество замкнутых экземпляров.

Итак, обе импликации доказаны.
\end{proof}

\subsection{Пример применения теоремы Эрбрана}

Рассмотрим множество дизъюнктов:
$$
S = \{P(a),\ \neg P(x)\}.
$$

Эрбранов универсум:
$$
HU = \{a\}.
$$

Множество замкнутых экземпляров:
$$
S^* = \{P(a),\ \neg P(a)\}.
$$

Конечное подмножество
$$
S_0 = \{P(a),\neg P(a)\}
$$
невыполнимо в логике высказываний. Следовательно, по теореме Эрбрана множество $S$ невыполнимо.

\section{Метод резолюций}

\subsection{Идея метода резолюций}

\begin{definition}
\textbf{Метод резолюций} --- метод автоматического доказательства теорем, основанный на приведении отрицания доказываемого утверждения к противоречию с помощью правила резолюции.
\end{definition}

Обычно используется схема доказательства от противного:
$$
\Gamma \models A
$$
тогда и только тогда, когда
$$
\Gamma \cup \{\neg A\}
$$
невыполнимо.

Задача метода резолюций:
$$
\Gamma \cup \{\neg A\} \vdash \square.
$$

\subsection{Правило резолюции в логике высказываний}

Пусть имеются две клаузы:
$$
C_1 = A \lor L_1 \lor \dots \lor L_m,
$$
$$
C_2 = \neg A \lor M_1 \lor \dots \lor M_n.
$$

Тогда из них можно вывести резольвенту:
$$
L_1 \lor \dots \lor L_m \lor M_1 \lor \dots \lor M_n.
$$

Схематически:
$$
\frac{A \lor C \qquad \neg A \lor D}{C \lor D}.
$$

\subsubsection{Пример резолюции в логике высказываний}

Пусть:
$$
C_1 = P \lor Q,
$$
$$
C_2 = \neg P \lor R.
$$

Тогда резольвента:
$$
Q \lor R.
$$

Если имеются:
$$
P,\quad \neg P,
$$
то:
$$
\frac{P \qquad \neg P}{\square}.
$$

\subsection{Корректность правила резолюции}

\begin{theorem}[Корректность резолюции]
Если клаузы $A \lor C$ и $\neg A \lor D$ истинны в некоторой интерпретации, то их резольвента $C \lor D$ также истинна в этой интерпретации.
\end{theorem}

\begin{proof}
Пусть интерпретация $\mathcal{I}$ удовлетворяет:
$$
A \lor C
$$
и
$$
\neg A \lor D.
$$

Рассмотрим два случая.

\textbf{Случай 1.} $A$ истинно.

Тогда $\neg A$ ложно. Поскольку $\neg A \lor D$ истинно, а $\neg A$ ложно, должно быть истинно $D$. Следовательно, истинно $C \lor D$.

\textbf{Случай 2.} $A$ ложно.

Тогда, поскольку $A \lor C$ истинно, должно быть истинно $C$. Следовательно, истинно $C \lor D$.

В обоих случаях резольвента $C \lor D$ истинна.
\end{proof}

\subsection{Резолюция в логике первого порядка}

В логике первого порядка правило резолюции использует унификацию.

Пусть имеются клаузы:
$$
C_1 = L \lor C,
$$
$$
C_2 = \neg M \lor D.
$$

Если литералы $L$ и $M$ унифицируются с наиболее общим унификатором $\theta$, то можно вывести:
$$
(C \lor D)\theta.
$$

Схема:
$$
\frac{L \lor C \qquad \neg M \lor D}{(C \lor D)\theta},
$$
где
$$
L\theta = M\theta.
$$

\subsubsection{Пример резолюции первого порядка}

Пусть:
$$
C_1 = P(x) \lor Q(x),
$$
$$
C_2 = \neg P(a) \lor R.
$$

Литералы $P(x)$ и $P(a)$ унифицируются:
$$
\theta = \{x \mapsto a\}.
$$

Резольвента:
$$
(Q(x) \lor R)\theta = Q(a) \lor R.
$$

\subsection{Приведение формулы к клаузальной форме}

Метод резолюций обычно применяется к формулам в конъюнктивной нормальной форме, то есть к множеству клауз.

\subsubsection{Алгоритм приведения к клаузальной форме}

Для формулы логики первого порядка выполняются шаги:

\begin{enumerate}
    \item \textbf{Устранить импликации и эквиваленции.}
    $$
    A \to B \equiv \neg A \lor B,
    $$
    $$
    A \leftrightarrow B \equiv (A \to B)\land(B \to A).
    $$

    \item \textbf{Внести отрицания внутрь.}
    Используются законы де Моргана:
    $$
    \neg(A\land B) \equiv \neg A \lor \neg B,
    $$
    $$
    \neg(A\lor B) \equiv \neg A \land \neg B,
    $$
    $$
    \neg \forall x A \equiv \exists x \neg A,
    $$
    $$
    \neg \exists x A \equiv \forall x \neg A.
    $$

    \item \textbf{Переименовать связанные переменные}, чтобы разные кванторы связывали разные имена.

    \item \textbf{Вынести кванторы в пренексную форму.}
    Получить формулу вида
    $$
    Q_1x_1\dots Q_nx_n\, B,
    $$
    где $B$ не содержит кванторов.

    \item \textbf{Сколемизировать формулу.}
    Экзистенциальные переменные заменяются сколемовскими функциями от предшествующих универсальных переменных.

    Например:
    $$
    \forall x \exists y\, P(x,y)
    $$
    превращается в
    $$
    \forall x\, P(x,f(x)).
    $$

    Если экзистенциальная переменная не зависит от универсальных, она заменяется сколемовской константой:
    $$
    \exists x\, P(x)
    \quad\leadsto\quad
    P(c).
    $$

    \item \textbf{Удалить универсальные кванторы.}
    Все оставшиеся переменные считаются неявно универсально квантифицированными.

    \item \textbf{Привести матрицу к конъюнктивной нормальной форме.}
    Используются дистрибутивные законы:
    $$
    A \lor (B\land C) \equiv (A\lor B)\land(A\lor C).
    $$

    \item \textbf{Представить результат как множество клауз.}
\end{enumerate}

\subsubsection{Пример приведения к клаузальной форме}

Пусть дана формула:
$$
\forall x\,(P(x) \to \exists y\, Q(x,y)).
$$

Шаг 1:
$$
\forall x\,(\neg P(x) \lor \exists y\, Q(x,y)).
$$

Шаг 2: вынесение кванторов:
$$
\forall x \exists y\,(\neg P(x) \lor Q(x,y)).
$$

Шаг 3: сколемизация:
$$
\forall x\,(\neg P(x) \lor Q(x,f(x))).
$$

Шаг 4: удаление универсального квантора:
$$
\neg P(x) \lor Q(x,f(x)).
$$

Клауза:
$$
\{\neg P(x), Q(x,f(x))\}.
$$

\subsection{Полнота метода резолюций}

\begin{theorem}[Полнота резолюции для логики высказываний]
Если конечное множество дизъюнктов логики высказываний невыполнимо, то из него можно вывести пустую клаузу методом резолюций.
\end{theorem}

\begin{proof}
Дадим доказательство через построение дерева поиска по истинностным оценкам.

Пусть $S$ --- конечное невыполнимое множество клауз над атомами
$$
p_1,\dots,p_n.
$$

Рассмотрим все возможные истинностные оценки этих атомов. Так как $S$ невыполнимо, каждая полная оценка делает ложной хотя бы одну клаузу из $S$.

Клауза
$$
L_1 \lor \dots \lor L_k
$$
ложна тогда и только тогда, когда все её литералы ложны. Поэтому каждой ветви полного дерева оценок соответствует клауза, запрещающая эту ветвь.

Покажем индукцией по числу атомов, что если множество клауз запрещает все оценки, то из него резолюцией выводится пустая клауза.

\textbf{База.} Если атомов нет, то единственная возможная клауза без литералов --- пустая клауза $\square$. Значит, она уже содержится в множестве или выводится немедленно.

\textbf{Переход.} Пусть утверждение верно для $n-1$ атома. Рассмотрим атом $p_n$.

Разобьём все оценки на две группы:
$$
p_n = true
\quad\text{и}\quad
p_n = false.
$$

Так  как множество $S$ невыполнимо, обе группы оценок запрещены.

Рассмотрим два множества клауз:
\[
S_{p_n=true}
\]
и
\[
S_{p_n=false},
\]
получающиеся из $S$ подстановкой соответственно $p_n=true$ и $p_n=false$ с упрощением клауз.

При подстановке $p_n=true$:
\begin{itemize}
    \item всякая клауза, содержащая $p_n$, становится истинной и удаляется;
    \item из клауз, содержащих $\neg p_n$, литерал $\neg p_n$ удаляется;
    \item остальные клаузы остаются без изменения.
\end{itemize}

Аналогично, при подстановке $p_n=false$:
\begin{itemize}
    \item всякая клауза, содержащая $\neg p_n$, становится истинной и удаляется;
    \item из клауз, содержащих $p_n$, литерал $p_n$ удаляется;
    \item остальные клаузы остаются без изменения.
\end{itemize}

Оба полученных множества невыполнимы над атомами
\[
p_1,\dots,p_{n-1}.
\]
По индукционному предположению из каждого из них выводится пустая клауза.

Теперь свяжем эти выводы с исходным множеством $S$.

Вывод пустой клаузы из $S_{p_n=true}$ соответствует выводу некоторой клаузы
\[
\neg p_n
\]
из исходного множества $S$.

Интуитивно это означает: если при предположении $p_n=true$ получается противоречие, то из $S$ следует $\neg p_n$.

Аналогично вывод пустой клаузы из $S_{p_n=false}$ соответствует выводу клаузы
\[
p_n
\]
из $S$.

После этого применяем резолюцию:
\[
\frac{p_n \qquad \neg p_n}{\square}.
\]

Следовательно, из $S$ выводится пустая клауза. Теорема доказана.
\end{proof}

\begin{theorem}[Полнота резолюции для логики первого порядка]
Если множество клауз логики первого порядка невыполнимо, то из него методом резолюций первого порядка можно вывести пустую клаузу.
\end{theorem}

\begin{proof}
Пусть $S$ --- невыполнимое множество клауз логики первого порядка.

По теореме Эрбрана существует конечное множество замкнутых экземпляров
\[
S_0 \subseteq S^*
\]
такое, что $S_0$ невыполнимо как множество клауз логики высказываний.

По полноте резолюции для логики высказываний из $S_0$ можно вывести пустую клаузу с помощью пропозициональной резолюции:
\[
S_0 \vdash_{res} \square.
\]

Каждая клауза из $S_0$ является замкнутым экземпляром некоторой клаузы из $S$. Пропозициональный резолюционный вывод над замкнутыми экземплярами можно поднять к выводу над исходными клаузами первого порядка, используя унификацию и наиболее общие унификаторы.

Это свойство называется \textbf{леммой подъёма}: если резольвента выводится из замкнутых экземпляров двух клауз, то существует резолюционный шаг над самими исходными клаузами, некоторый замкнутый экземпляр которого совпадает с данной пропозициональной резольвентой.

Применяя лемму подъёма ко всем шагам вывода из $S_0$, получаем вывод пустой клаузы из $S$:
\[
S \vdash_{res} \square.
\]

Следовательно, резолюция первого порядка полна.
\end{proof}

\subsection{Алгоритм доказательства методом резолюций}

Пусть требуется доказать:
\[
\Gamma \models A.
\]

Алгоритм:
\begin{enumerate}
    \item Сформировать множество формул:
    \[
    \Gamma \cup \{\neg A\}.
    \]
    \item Привести все формулы к клаузальной форме.
    \item Получить множество клауз $S$.
    \item Пока возможно:
    \begin{enumerate}
        \item выбрать две клаузы $C_1$ и $C_2$;
        \item выбрать пару контрарных литералов;
        \item найти наиболее общий унификатор $\theta$;
        \item построить резольвенту;
        \item добавить её к множеству клауз, если она новая;
        \item если получена пустая клауза $\square$, завершить доказательство успешно.
    \end{enumerate}
    \item Если больше нельзя получить новых клауз и пустая клауза не получена, доказательство не найдено.
\end{enumerate}

\subsection{Пример доказательства методом резолюций}

Докажем утверждение:

\[
\forall x\,(Man(x) \to Mortal(x)), \quad Man(socrates)
\models Mortal(socrates).
\]

\subsubsection{Шаг 1. Добавляем отрицание цели}

Исходные предпосылки:
\[
\forall x\,(Man(x) \to Mortal(x)),
\]
\[
Man(socrates).
\]

Отрицание цели:
\[
\neg Mortal(socrates).
\]

\subsubsection{Шаг 2. Приведение к клаузам}

Первая формула:
\[
\forall x\,(\neg Man(x) \lor Mortal(x)).
\]

Клауза:
\[
C_1 = \neg Man(x) \lor Mortal(x).
\]

Вторая формула:
\[
C_2 = Man(socrates).
\]

Отрицание цели:
\[
C_3 = \neg Mortal(socrates).
\]

\subsubsection{Шаг 3. Резолюционный вывод}

Берём:
\[
C_1 = \neg Man(x) \lor Mortal(x)
\]
и
\[
C_2 = Man(socrates).
\]

Контрарные литералы:
\[
\neg Man(x)
\quad\text{и}\quad
Man(socrates).
\]

Унификатор:
\[
\theta_1 = \{x \mapsto socrates\}.
\]

Резольвента:
\[
C_4 = Mortal(socrates).
\]

Теперь берём:
\[
C_4 = Mortal(socrates)
\]
и
\[
C_3 = \neg Mortal(socrates).
\]

Получаем:
\[
\square.
\]

Следовательно:
\[
\forall x\,(Man(x) \to Mortal(x)), \quad Man(socrates)
\models Mortal(socrates).
\]

\section{Логическое программирование и резолюция в Prolog}

\subsection{Связь Prolog с хорновской резолюцией}

Prolog-программы являются, в основном, множествами определённых хорновских клауз вида:
\[
A \leftarrow B_1,\dots,B_n.
\]

Запрос:
\[
?- G_1,\dots,G_k
\]
соответствует отрицательной клаузе:
\[
\leftarrow G_1,\dots,G_k.
\]

Вычисление в Prolog можно понимать как построение SLD-вывода:
\[
P \cup \{\leftarrow G\} \vdash \square.
\]

Если вывод успешен, то найденная композиция подстановок даёт ответ на запрос.

\subsection{Backtracking}

\begin{definition}
\textbf{Backtracking}, или возврат, --- механизм перебора альтернатив, при котором система возвращается к предыдущей точке выбора, если текущий путь поиска не приводит к успеху.
\end{definition}

\subsubsection{Пример backtracking}

Программа:
\begin{lstlisting}[language=Prolog]
color(red).
color(green).
color(blue).
\end{lstlisting}

Запрос:
\begin{lstlisting}[language=Prolog]
?- color(X).
\end{lstlisting}

Ход поиска:
\[
X = red.
\]

При запросе следующего решения Prolog возвращается к точке выбора и пробует следующий факт:
\[
X = green.
\]

Затем:
\[
X = blue.
\]

После этого альтернатив больше нет.

\subsection{Рекурсия в Prolog}

В Prolog циклы обычно выражаются через рекурсию.

Пример определения принадлежности элемента списку:
\begin{lstlisting}[language=Prolog]
member(X, [X|_]).
member(X, [_|Tail]) :- member(X, Tail).
\end{lstlisting}

Запрос:
\begin{lstlisting}[language=Prolog]
?- member(2, [1,2,3]).
\end{lstlisting}

Шаги:
\[
member(2,[1,2,3])
\]
не унифицируется с первым правилом, потому что $2 \neq 1$.

Применяется второе правило:
\[
member(2,[1,2,3]) \Rightarrow member(2,[2,3]).
\]

Теперь:
\[
member(2,[2,3])
\]
унифицируется с первым правилом:
\[
member(X,[X|_]).
\]

Следовательно, запрос успешен.

\subsection{Отсечение}

\begin{definition}
\textbf{Отсечение} \texttt{!} в Prolog --- управляющий оператор, который запрещает возврат к альтернативам, созданным до него в текущем предикате.
\end{definition}

Пример:
\begin{lstlisting}[language=Prolog]
max(X, Y, X) :- X >= Y, !.
max(_, Y, Y).
\end{lstlisting}

Если условие $X \geq Y$ выполнено, Prolog фиксирует первый вариант и не рассматривает второе правило.

\subsection{Отрицание как неудача}

В Prolog отрицание часто реализуется как \textbf{отрицание как неудача}:
\[
not(G)
\]
истинно, если попытка доказать $G$ завершается неудачей.

Это не классическое логическое отрицание. Оно зависит от полноты базы знаний и от стратегии поиска.

Пример:
\begin{lstlisting}[language=Prolog]
bird(tweety).
penguin(polly).

flies(X) :- bird(X), not(penguin(X)).
\end{lstlisting}

Запрос:
\begin{lstlisting}[language=Prolog]
?- flies(tweety).
\end{lstlisting}
успешен, если невозможно доказать:
\begin{lstlisting}[language=Prolog]
penguin(tweety).
\end{lstlisting}

\section{Сравнение языков и парадигм программирования}

\subsection{Сравнение рассматриваемых языков}

\begin{center}
\begin{tabular}{|m{2.5cm}|m{3cm}|m{6cm}|}
\hline
\textbf{Язык} & \textbf{Парадигма} & \textbf{Ключевые особенности} \\
\hline
Fortran & Процедурная & Численные вычисления, массивы, высокая производительность \\
\hline
C & Процедурная, системная & Указатели, ручное управление памятью, близость к аппаратуре \\
\hline
LISP & Функциональная & Списки, S-выражения, рекурсия, макросы \\
\hline
ML & Функциональная & Статическая типизация, вывод типов, сопоставление с образцом \\
\hline
Haskell & Чистая функциональная & Ленивые вычисления, монады, классы типов \\
\hline
Prolog & Логическая & Факты, правила, унификация, SLD-резолюция, backtracking \\
\hline
Java & Объектно-ориентированная & JVM, байт-код, сборка мусора, интерфейсы \\
\hline
C\# & Объектно-ориентированная & CLR, .NET, свойства, делегаты, LINQ, async/await \\
\hline
\end{tabular}
\end{center}

\subsection{Сравнение способов описания факториала}

\subsubsection{C}

\begin{lstlisting}[language=C]
int fact(int n) {
    int r = 1;
    for (int i = 2; i <= n; i++) {
        r *= i;
    }
    return r;
}
\end{lstlisting}

Здесь вычисление описано как последовательное изменение переменной \texttt{r}.

\subsubsection{Haskell}

\begin{lstlisting}[language=Haskell]
fact 0 = 1
fact n = n * fact (n - 1)
\end{lstlisting}

Здесь вычисление задано рекурсивным математическим определением.

\subsubsection{Prolog}

\begin{lstlisting}[language=Prolog]
fact(0, 1).
fact(N, F) :-
    N > 0,
    N1 is N - 1,
    fact(N1, F1),
    F is N * F1.
\end{lstlisting}

Запрос:
\begin{lstlisting}[language=Prolog]
?- fact(4, F).
\end{lstlisting}

Результат:
\[
F = 24.
\]

\subsection{Выбор парадигмы}

Выбор языка и парадигмы зависит от характера задачи.

\begin{itemize}
    \item Для численного моделирования подходят Fortran, C, C++.
    \item Для системного программирования подходит C.
    \item Для символьных вычислений и метапрограммирования подходят LISP-подобные языки.
    \item Для задач с большим количеством логических отношений подходит Prolog.
    \item Для больших прикладных систем часто применяются Java и C\#.
    \item Для задач, где важны математическая строгость, композиционность и работа с неизменяемыми структурами, удобны ML и Haskell.
\end{itemize}

\section{Краткие выводы}

\begin{enumerate}
    \item Процедурные языки описывают вычисление как последовательность команд, изменяющих состояние.
    \item Fortran исторически ориентирован на научные и численные вычисления.
    \item C является низкоуровневым процедурным языком, важным для системного программирования.
    \item Функциональные языки рассматривают программу как совокупность функций и выражений.
    \item LISP важен для символьной обработки и метапрограммирования.
    \item ML и Haskell демонстрируют развитые системы типов и функциональный стиль.
    \item Логическое программирование задаёт программу как систему фактов и правил.
    \item Prolog использует унификацию, SLD-резолюцию и backtracking.
    \item Объектно-ориентированное программирование строит программы из объектов, объединяющих состояние и поведение.
    \item Java и C\# являются распространёнными объектно-ориентированными языками с управляемой средой выполнения.
    \item Теорема Эрбрана связывает невыполнимость множества дизъюнктов первого порядка с невыполнимостью конечного множества их замкнутых экземпляров.
    \item Метод резолюций является основой многих систем автоматического доказательства теорем.
\end{enumerate}

\end{document}
